\section{Method}
\label{sec:method}

\subsection{Overview}

MemGen generates latent memories through a Trigger--Weaver--Reasoner pathway \cite{zhang2025memgen}. During inference, the Trigger monitors the current reasoning state and determines whether latent-memory augmentation should be invoked. When augmentation is enabled, the Weaver transforms the current hidden state into a latent memory, which is subsequently injected into the Reasoner.

In the original inference pathway, a generated latent memory is transient and is not explicitly retained for later reasoning steps. We extend this pathway with session-local latent memory bank. The bank stores Weaver-generated latent memories from earlier context-processing steps, retrieves relevant memories according to the current reasoning state, and feeds the retrieved memories back into the Weaver. The Weaver therefore conditions on both the current hidden state and previously generated latent memories before producing a new latent memory for the Reasoner.

The resulting method changes the inference flow without modifying the parameters of the Trigger, Weaver, or Reasoner. It introduces no parameter updates, additional training objective, or cross-session memory. All persistent memory state is confined to the current inference session and is reset before processing an independent context.

\subsection{Session-Local Latent Memory Bank}

For an input context, let $H_t$ denote the Reasoner's hidden state at inference step $t$. When the Trigger emits an invocation signal, the memory bank is queried using a representation derived from the recent reasoning state.

The query representation is obtained by mean-pooling the most recent $L$ hidden states:

\[
q_t = \operatorname{MeanPool}\left(H_t[-L:]\right).
\]

The memory bank at step $t$ is denoted by

\[
\mathcal{M}_t =
\left\{
\left(m_i,k_i,r_i^{\mathrm{last}}\right)
\right\}_{i=1}^{N_t},
\]

where $m_i$ is a detached Weaver-space latent memory, $k_i$ is its retrieval key, and $r_i^{\mathrm{last}}$ records the global retrieval count at which memory $m_i$ was most recently retrieved.

Each retrieval key is obtained by mean-pooling the latent-token vectors of the corresponding memory:

\[
k_i = \operatorname{MeanPool}(m_i).
\]

Before storage, each generated latent memory is detached from the computation graph and transferred to memory:

\[
m_i^{\mathrm{store}} =
{Detach}(m_i)
\]

The bank belongs to exactly one context or inference session. It has no global registry, does not share slots across contexts, and is reset before the next independent sample. Its capacity is bounded by a configurable maximum number of slots $C$. In the evaluated configuration, $C=16$, and each slot stores one Weaver-generated memory consisting of eight latent-token vectors.

\subsection{Similarity-Based Retrieval with Temporal Decay}

For each stored memory, the bank computes a retrieval score by combining cosine similarity with temporal decay:

\[
s_i(q_t) =
\operatorname{cosine}(q_t,k_i)
\exp\left(-\alpha \Delta r_i\right),
\]

where

\[
\Delta r_i =
r_t-r_i^{\mathrm{last}}.
\]

Here, $r_t$ denotes the current global retrieval count, and $r_i^{\mathrm{last}}$ denotes the retrieval count at which memory $m_i$ was most recently selected. Accordingly, $\Delta r_i$ measures the retrieval inactivity of memory $m_i$, namely, the number of retrieval operations that have occurred since it was last retrieved. The term $\exp(-\alpha \Delta r_i)$ represents the corresponding temporal-decay weight: memories that have remained unretrieved for longer receive smaller weights. The coefficient $\alpha$ controls the strength of this decay.

The retrieval procedure first filters memories using a retrieval threshold $\tau_r$:

\[
\mathcal{C}_t =
\left\{
m_i\in\mathcal{M}_t
\mid
s_i(q_t)\geq\tau_r
\right\}.
\]

Among the surviving candidates, the bank returns at most the $K$ highest-scoring memories:

\[
R_t =
\operatorname{TopK}\left(\mathcal{C}_t,K\right).
\]

If no memory passes the retrieval threshold, the bank returns an empty retrieval set:

\[
R_t=\varnothing.
\]

The method does not force a top-1 result when every stored memory falls below the threshold. Whenever a memory is successfully retrieved, its last-retrieval count is updated to the current global retrieval count:

\[
r_i^{\mathrm{last}}\leftarrow r_t,
\qquad
m_i\in R_t.
\]

In the evaluated configuration, we set

\[
\tau_r=0.05,
\qquad
K=2,
\qquad
\alpha=0.05.
\]

These values remain fixed across all evaluated EventQA contexts.

\subsection{Retrieved-Memory-Conditioned Latent Generation}

The central difference from the original transient MemGen pathway is that retrieved memories are returned to the Weaver rather than being directly injected into the Reasoner.

At inference step $t$, the Trigger first evaluates the current reasoning state:

\[
g_t=\operatorname{Trigger}(H_t),
\]

where $g_t\in\{\texttt{INVOKE},\texttt{SKIP}\}$.

When the Trigger emits \texttt{INVOKE}, the memory bank retrieves $R_t$ using the current reasoning-state query. The Weaver then conditions on both the current hidden state and the retrieved latent memories:

\[
m_t =
\operatorname{Weaver}(H_t,R_t).
\]

If no prior memory passes the retrieval threshold, the Weaver follows the original generation path:

\[
m_t =
\operatorname{Weaver}(H_t,\varnothing).
\]

The Weaver consolidates the current reasoning state and the retrieved latent memories into a newly generated latent memory $m_t$. This new latent memory is then passed through the existing Weaver-to-Reasoner pathway:

\[
y_t =
\operatorname{Reasoner}(H_t,m_t).
\]

When the Trigger emits \texttt{SKIP}, no latent-memory augmentation is performed:

\[
y_t =
\operatorname{Reasoner}(H_t).
\]

This design avoids directly injecting all retrieved memories into the Reasoner. Instead, the frozen Weaver transforms the current hidden state together with the retrieved latent support into a new latent memory. This intermediate Weaver transformation integrates the retrieved memories before they are used by the Reasoner, reducing potential interference while preserving the original MemGen latent-generation pathway.

The method therefore reuses generated latent-token memories without decoding them into natural-language text or exposing their original source context again.

\subsection{Memory Write and Thread Update}

A newly generated latent memory $m_t$ is written to the bank whenever the Weaver is invoked. The write policy determines whether the candidate represents a new latent-memory thread or should refresh an existing one.

The key of the new memory is

\[
k_t =
\operatorname{MeanPool}(m_t).
\]

If the bank is empty, the new memory is inserted directly:

\[
\mathcal{M}_{t+1} =
\left\{
\left(m_t,k_t,r_t\right)
\right\}.
\]

Otherwise, the bank identifies the stored memory with the greatest similarity to the new candidate:

\[
j =
\arg\max_i
\operatorname{cosine}(k_t,k_i).
\]

Let

\[
u_t =
\max_i
\operatorname{cosine}(k_t,k_i)
\]

denote the maximum candidate-to-memory similarity. The bank uses a update threshold $\tau_u$.

If the maximum similarity is lower than the update threshold, the candidate is treated as a distinct latent-memory thread and inserted as a new slot:

\[
u_t<\tau_u
\quad\Rightarrow\quad
\mathcal{M}_{t+1} =
\mathcal{M}_t
\cup
\left\{
\left(m_t,k_t,r_t\right)
\right\}.
\]

If the maximum similarity reaches or exceeds the update threshold, the most similar slot is refreshed by replacing its stored latent memory:

\[
u_t\geq\tau_u
\quad\Rightarrow\quad
\mathcal{M}_{t+1} =
\left(
\mathcal{M}_t
\setminus
\left\{
\left(m_j,k_j,r_j^{\mathrm{last}}\right)
\right\}
\right)
\cup
\left\{
\left(m_t,k_t,r_t\right)
\right\}.
\]

Thus, sufficiently dissimilar generated memories create new slots, whereas similar generated memories replace and refresh an existing latent-memory thread. In the evaluated configuration, the update threshold is fixed to

\[
\tau_u=0.10.
\]

The retrieval threshold $\tau_r$ and update threshold $\tau_u$ serve different purposes. The retrieval threshold determines whether an existing memory is sufficiently relevant to the current reasoning state, whereas the update threshold determines whether a newly generated memory should replace an existing thread or create a new one.

\subsection{Bounded-Capacity Replacement}

The memory bank maintains a fixed capacity $C$:

\[
|\mathcal{M}_t| \leq C.
\]

When inserting a new memory thread would exceed this capacity, the bank first selects an existing slot for eviction. For each slot $i$, we define its retrieval age as

\[
a_i = r_t - r_i^{\mathrm{last}},
\]

where $r_t$ is the current global retrieval count and $r_i^{\mathrm{last}}$ is updated only when slot $i$ is selected during retrieval. A larger $a_i$ therefore indicates that the memory has remained unretrieved for more retrieval operations.

The bank first identifies the maximum retrieval age:

\[
a_{\max} = \max_i a_i.
\]

The newly generated memory then replaces the selected slot:

\[
(m_e,k_e,r_e^{\mathrm{last}})
\leftarrow
(m_t,k_t,r_t).
\]

This deterministic policy preferentially preserves recently retrieved latent memories rather than simply retaining the most recently written slots.

\subsection{End-to-End Inference Procedure}

The complete inference procedure is summarized below.

\begin{quote}
\ttfamily
Input:\\
\quad hidden states \(H_1, \ldots, H_T\)\\
\quad retrieval threshold \(\tau_r\)\\
\quad update threshold \(\tau_u\)\\
\quad temporal-decay coefficient \(\alpha\)\\
\quad maximum retrieval count \(K\)\\
\quad memory capacity \(C\)\\
Initialize:\\
\quad \(M \leftarrow \varnothing\)\\
\quad \(\texttt{retrieval\_count} \leftarrow 0\)\\
for \(t = 1, \ldots, T\):\\
\quad \(\texttt{signal} \leftarrow \operatorname{Trigger}(H_t)\)\\
\quad if \texttt{signal == INVOKE}:\\
\qquad \(q_t \leftarrow \operatorname{MeanPool}(H_t[-L:])\)\\
\qquad \(\texttt{retrieval\_count} \leftarrow \texttt{retrieval\_count} + 1\)\\
\qquad for each memory \(m_i\) in \(M\):\\
\qquad\quad \(k_i \leftarrow \operatorname{MeanPool}(m_i)\)\\
\qquad\quad \(\Delta r_i \leftarrow \texttt{retrieval\_count} - m_i.\texttt{last\_retrieved}\)\\
\qquad\quad \(s_i \leftarrow \operatorname{cosine}(q_t, k_i) \cdot \exp(-\alpha\Delta r_i)\)\\
\qquad \(R_t \leftarrow \operatorname{TopK}(\{m_i : s_i \geq \tau_r\}, K)\)\\
\qquad for each retrieved memory \(m_i\) in \(R_t\):\\
\qquad\quad \(m_i.\texttt{last\_retrieved} \leftarrow \texttt{retrieval\_count}\)\\
\qquad \(m_t \leftarrow \operatorname{Weaver}(H_t, R_t)\)\\
\qquad \(y_t \leftarrow \operatorname{Reasoner}(H_t, m_t)\)\\
\qquad \(k_t \leftarrow \operatorname{MeanPool}(m_t)\)\\
\qquad if \(M\) is empty:\\
\qquad\quad insert \(\operatorname{detach}(m_t).\texttt{cpu()}\) into \(M\)\\
\qquad else:\\
\qquad\quad \(j \leftarrow \arg\max_i \operatorname{cosine}(k_t, k_i)\)\\
\qquad\quad \(u_t \leftarrow \operatorname{cosine}(k_t, k_j)\)\\
\qquad\quad if \(u_t \geq \tau_u\):\\
\qquad\qquad replace \(m_j\) with \(\operatorname{detach}(m_t).\texttt{cpu()}\)\\
\qquad\quad else:\\
\qquad\qquad if \(|M| = C\):\\
\qquad\qquad\quad evict one slot using the bounded replacement policy\\
\qquad\qquad insert \(\operatorname{detach}(m_t).\texttt{cpu()}\) into \(M\)\\
\quad else:\\
\qquad \(y_t \leftarrow \operatorname{Reasoner}(H_t)\)
\end{quote}
